% LaTeX file for resume
\documentclass[sectioned]{dsyangres}
\usepackage[left =1in, top = 0.4in, bottom=0.4in]{geometry}

% Tectonic uses the XeTeX engine, so load the real TeX Gyre Pagella OTF
% via fontspec (same font Overleaf embeds). The classic Type1 tgpagella
% package silently falls back to Latin Modern under XeTeX, which looks
% worse -- so fontspec is required here despite the extra dependencies.
\usepackage{fontspec}
\usepackage{array}
\setmainfont
[ Mapping = tex-text ,
  Numbers = OldStyle ,
  Ligatures = {Common,Contextual} ,
  BoldFont = texgyrepagella-bold.otf ,
  ItalicFont = texgyrepagella-italic.otf ,
  BoldItalicFont = texgyrepagella-bolditalic.otf ]
{texgyrepagella-regular.otf}

% the margin option causes section titles to appear to the left of body text
% Sized so the name/contact header (which spans \textwidth plus the
% 0.95in section-title gutter) lands flush with a symmetric 1in right
% margin, matching the geometry package's 1in left margin.
\textwidth=5.55in

% Fixed-width title/date columns for job & project entries, so the date
% always lines up with the title's first line (instead of a free-floating
% \hfill) and bullet text never runs into the date column. The tabular
% spans the full \textwidth (title + gap + date, no extra slack) so the
% date column's right edge lands on the same symmetric right margin as
% the header above it.
\newlength{\jobdateswidth}
\setlength{\jobdateswidth}{0.9in}
\newlength{\jobheadergap}
\setlength{\jobheadergap}{0.3in}
\newlength{\jobtitlewidth}
\setlength{\jobtitlewidth}{\textwidth}
\addtolength{\jobtitlewidth}{-\jobdateswidth}
\addtolength{\jobtitlewidth}{-\jobheadergap}

\newcommand{\jobheader}[2]{%
  \noindent
  \begin{tabular}[t]{@{}>{\raggedright\arraybackslash\hyphenpenalty=10000\exhyphenpenalty=10000}p{\jobtitlewidth}@{\hspace{\jobheadergap}}>{\raggedleft\arraybackslash}p{\jobdateswidth}@{}}
  #1 & \textbf{#2}
  \end{tabular}\par
}

% Wraps a job/project entry: prints the title/date header row, then narrows
% \hsize/\linewidth/\textwidth so the following itemize wraps within the
% same column as the title (not a minipage, so entries can still break
% across a page if needed). Hyphenation is disabled so words wrap whole
% instead of splitting across lines.
\newenvironment{jobentry}[2]{%
  \jobheader{#1}{#2}%
  \begingroup
  \hsize=\jobtitlewidth
  \linewidth=\jobtitlewidth
  \textwidth=\jobtitlewidth
  \tolerance=9999
  \emergencystretch=3em
  \hyphenpenalty=10000
  \exhyphenpenalty=10000
}{%
  \endgroup
}

\begin{document}

\name{Daniel S. Yang}
\email{dsyang92@gmail.com}
\github{https://github.com/dsyang}
\linkedin{https://linkedin.com/in/dsyang92}
\phone{734-408-1407}


\begin{resume}

%%\section{Objective} Work on interesting problems with amazing people.

\section{Education}

\begin{jobentry}{\textbf{Carnegie Mellon University, Pittsburgh PA}}{Spring 2014}
  \begin{itemize}
    \item B.S. in Computer Science, Robotics minor
  \end{itemize}
\end{jobentry}

\section{Work \\ Experience}
\begin{jobentry}{\textbf{Notion}: Software Engineer - Mobile, San Francisco, CA}{2026 - Present}
  \begin{itemize}
      \item Shipped push notifications for agent messages on iOS Notion Core and Notion Agents app.
      \item Refactored search tab on Android to use modern Circuit presenter patterns.
      \item Improved bug reporting accuracy by embedding session ID into Android app screenshots.
  \end{itemize}
\end{jobentry}

\begin{jobentry}{\textbf{Notion}: Engineering Manager - Mobile Product, San Francisco, CA}{2024 - 2026}
  \begin{itemize}
      \item Defined mobile strategy and product direction with input from tech leads and xfn partners.
      \item Shifted team direction in 2025 to support Notion Agents natively, doubling mobile AI users and usage.
      \item Supported iOS and Android engineers across NY and SF. Grew team from 4 to 10.
  \end{itemize}
\end{jobentry}

\begin{jobentry}{\textbf{Notion}: Software Engineer - Android, San Francisco, CA}{2022 - 2024}
  \begin{itemize}
      \item Productionized Notion's mobile-native strategy by optimizing Jetpack Compose performance for native home tab.
      \item Improved app start time from $\textgreater$ 10 seconds to $\textless$ 2 seconds.
      \item Led iOS \& Android engineers to rebuild Notifications tab. Built an API endpoint for logic-sharing and future extensibility.
  \end{itemize}
\end{jobentry}

\begin{jobentry}{\textbf{Instagram}: Software Engineer, San Francisco, CA}{2019 - 2022}
  \begin{itemize}
      \item Tech led Android client infrastructure projects for Direct Messaging to support cross-app communication and multiple backends.
  \end{itemize}
\end{jobentry}

\begin{jobentry}{\textbf{Facebook}: Engineer in Residence, Seaside, CA}{2019}
  \begin{itemize}
      \item Taught Advanced Algorithms to over 100 students as a visiting instructor at California State University -- Monterey Bay.
  \end{itemize}
\end{jobentry}

\begin{jobentry}{\textbf{Facebook}: Software Engineer, Menlo Park, CA}{2014 - 2018}
  \begin{itemize}
      \item Built multiple features for photos and sharing on Android.
      \item Drove dramatic scroll performance improvements in the Android app through the adoption of Litho across dozens of teams.
  \end{itemize}
\end{jobentry}


\section{Hobby \\ Projects}

\begin{jobentry}{\textbf{code.imdaniel.fyi} (https://code.imdaniel.fyi)}{2025}
  \begin{itemize}
    \item Site of small, static HTML utilities and games, with iOS and Android apps to use them offline.
    \item An exercise in building and shipping products entirely from a phone using the Claude and GitHub mobile apps.
  \end{itemize}
\end{jobentry}

\begin{jobentry}{\textbf{readcode} (https://github.com/dsyang/readcode)}{2025}
  \begin{itemize}
    \item Code review tool for commenting across multiple commits and making manual edits. Built to review and tweak AI-generated code without an IDE.
    \item Rust + Tauri desktop app for Windows and macOS with auto-update and in-app feedback reporting.
  \end{itemize}
\end{jobentry}

\begin{jobentry}{\textbf{Buck} (https://buck.build)}{2016}
  \begin{itemize}
    \item Added original build rules for Kotlin, unblocking multiple companies such as Uber to use buck.
  \end{itemize}
\end{jobentry}

\section{Skills}
\textbf{Experienced} with Android, iOS, Java, Kotlin, Jetpack Compose, Circuit, \\ Swift, TCA, Tuist, Typescript, React, Rust, Standard ML, Python.

\end{resume}
\end{document}
